\chapter{绪论}

\section{研究背景与意义}

在社交媒体时代，信息的传播并非简单的广播式扩散，而是一个受用户认知过程深度影响的动态行为\cite{hu2017survey}。用户在接收到一条评论或帖子后，需要经历注意捕获、内容理解、态度形成与判断等认知加工环节，才可能进一步产生回复、转发等传播行为。这一由认知加工驱动的传播过程被称为信息认知传播，它构成了信息在社会网络中扩散的核心机制。对信息认知传播特征的量化理解，有助于揭示社交媒体上信息扩散的内在规律，对舆情监测、内容推荐和平台治理等应用场景具有重要的理论与实践价值。

近年来，大语言模型技术的快速发展催生了新型的信息传播主体\cite{brown2020language}。大语言模型驱动的AI代理已在多个社交媒体平台上具备独立发布评论与参与讨论的能力（如X平台上的Grok，见图\ref{fig:x_grok_example}），信息传播主体的构成由此发生结构性变化——人类用户与AI代理作为具有不同认知与内容生成机制的信息发布主体共存于同一传播系统中。

\begin{figure}[htbp]
  \centering
  \includegraphics[width=0.55\textwidth]{G1-0/X-grok.png}
  \caption{社交媒体平台上人类用户与AI代理共存的讨论实例}
  \label{fig:x_grok_example}
\end{figure}

人类用户受社交动机、情绪状态与认知负荷驱动，AI代理则以确定性的语言模型生成机制参与讨论\cite{brown2020language}，两类主体不同的信息处理机制会产生差异化的传播行为模式。已有研究表明，社交机器人与人类用户在发布频率、内容模式与网络结构等行为特征上呈现显著差异\cite{ferrara2016rise,himelein2021bots}，大语言模型驱动的新一代AI代理在内容生成能力上已接近人类水平，使得传统的人机行为边界日趋模糊\cite{yang2024anatomy}。此前针对传统社交机器人的研究已发现，自动化账号在信息传播早期通过大量转发低可信度来源的内容并定向@有影响力的人类用户，充当了虚假信息的放大器\cite{shao2018spread}，表明非人类主体的参与会对信息传播生态产生实质性影响。这一多主体共存的格局为信息认知传播研究提出了新的问题：通过对比人类用户与AI代理在传播行为指标上的系统性差异，可以从多主体视角揭示信息认知传播的新特征。

图\ref{fig:comment_tree_comparison}以示意图的方式说明了同质主体与多主体共存两种场景下评论树结构的差异。左侧为传统的同质讨论树，所有节点均由人类用户发布，节点间的交互模式单一（仅存在 H$\to$H）；右侧为多主体共存的异构讨论树，人类用户（H）与 AI 代理（A）共同参与讨论，路径上出现 H$\to$H、H$\to$A、A$\to$H 等多种交互类型（由于 AI 代理主动回复另一 AI 代理的情形在当前平台中较为罕见，图中未包含 A$\to$A 交互），传播行为模式因主体类型及其组合而产生分化。这一结构性变化构成了本文研究的出发点。

\begin{figure}[htbp]
  \centering
  \includegraphics[width=0.9\textwidth]{G1-1/图1-1 同质主体与多主体共存场景下的评论树结构对比.cropped.pdf}
  \caption{同质主体与多主体共存场景下的评论树结构示意图}
  \label{fig:comment_tree_comparison}
\end{figure}

从图\ref{fig:comment_tree_comparison}的异构讨论树中可以观察到两个层次的分析维度。在节点层面，每条评论由特定类型的主体发布，不同主体生成的评论在内容质量、表达风格和引发后续回复的能力上可能存在差异，这要求在评估单条评论的传播价值时将发布主体类型纳入考量。在路径层面，从根节点到任意节点的讨论链上，相邻节点间的回复关系因主体类型的组合而呈现异构性——如 H$\to$A 边表示人类回复AI生成的内容，A$\to$H 边表示AI回复人类发起的讨论——不同交互类型对讨论能否继续延续可能具有不同的影响。然而，现有的信息传播分析与级联预测方法普遍假设传播节点具有同质性，未将主体类型纳入建模框架，在消息层评估、讨论链预测和数据支撑三个层面均存在不足（详见1.2节综述）。

针对上述不足，本文设计并实现了多主体信息认知传播评估系统。由于受众的认知加工过程发生于个体内部、难以直接观测，本文采用传播行为指标作为认知过程的代理变量——回复、讨论延续等可观测行为是认知加工的外化结果，通过分析这些指标在不同主体间的分布差异，可以间接刻画多主体环境下的信息认知传播特征。基于这一思路，本文从消息层、级联层和交互类型层三个评估维度对人类用户与AI代理的传播行为差异进行量化评估。本研究的理论意义在于：为多主体共存的社交媒体环境提供了一套系统性的传播行为评估工具，以传播行为指标间接刻画多主体生态下的信息认知传播特征，填补了现有分析方法在多主体场景下的能力空白。本研究的社会意义在于：在AI代理日益深入参与社交媒体讨论的背景下，量化AI代理与人类用户的传播行为差异有助于平台治理方识别AI驱动的信息扩散模式，为制定多主体环境下的内容治理策略提供数据支撑；同时，揭示人机共存对信息传播生态的影响，为评估AI代理参与公共讨论的社会效应提供实证依据，对当前关于AI生成内容标注与监管的政策讨论具有参考价值。

\section{国内外研究现状}

本节从信息传播与级联预测、社交媒体多主体行为分析两个方向综述国内外相关研究，梳理现有工作的进展与不足，明确本文的研究切入点。

\subsection{信息传播与级联预测研究现状}

信息传播建模是社交媒体研究的核心议题之一。胡长军等人\cite{hu2017survey}系统梳理了在线社交网络信息传播研究的主要方向，指出信息传播过程受网络结构、用户行为与内容特征的共同驱动。Goel等人\cite{goel2016structural}提出"结构病毒性"（Structural Virality）概念，将传播结构视为广播式与病毒式传播之间的连续谱，基于大规模Twitter级联的实证分析发现绝大多数传播事件由广播式扩散驱动，为理解信息扩散的宏观模式提供了定量框架。

在级联预测方向，研究者探索了多种深度学习方法。Li等人\cite{li2017deepcas}提出DeepCas，通过随机游走与注意力GRU网络首次实现级联规模的端到端深度学习预测，显著优于基于手工特征的传统方法；Cao等人\cite{cao2017deephawkes}提出DeepHawkes，在Hawkes过程框架下融合深度学习建模用户影响力与传播路径的时序依赖，弥合了统计模型与深度学习之间的差距；Zhou等人\cite{zhou2021casflow}提出CasFlow，引入变分自编码器框架显式捕捉级联树状拓扑与传播不确定性，在多个数据集上取得领先性能。

传播结构分析在可信度判别领域同样发挥了重要作用。Ma等人\cite{ma2017detect}提出传播树核方法并构建了Twitter15/16基准数据集，通过度量传播树的结构相似性实现谣言检测；Ma等人\cite{ma2016detecting}利用循环神经网络建模微博传播中用户响应的时序动态，实现早期谣言检测。在评论级交互分析方面，Hsu等人\cite{hsu2009ranking}提出结合文本特征、用户声誉与社区反馈的评论排序方法；Glenski和Weninger\cite{glenski2017predicting}分析了用户、内容与社区上下文特征对评论互动量的影响，发现社区动态因素具有不可忽视的作用。

然而，上述方法普遍存在同质性假设的局限：无论是级联预测模型还是评论分析方法，均将传播网络中的参与节点视为同质主体，未考虑不同类型主体在内容生成机制和传播行为模式上的结构性差异，在传播路径建模中忽略了相邻节点主体类型组合对传播延续的差异化影响。

\subsection{社交媒体多主体行为分析研究现状}

社交机器人的兴起是社交媒体从同质主体向多主体共存环境演变的关键驱动因素。Ferrara等人\cite{ferrara2016rise}系统梳理了社交机器人的发展历程与社会影响，指出自动化账号已广泛参与政治讨论与舆论操纵，并归纳了其异常发布频率、重复内容模式等典型行为特征。Himelein-Wachowiak等人\cite{himelein2021bots}开展跨平台大规模对比研究，发现机器人在行为特征上表现出可区分的系统性模式，但差异因平台和地区不同而呈现显著变异。Shao等人\cite{shao2018spread}通过对1400万条Twitter消息的分析，揭示社交机器人通过大量转发低可信度内容并定向@有影响力的人类用户，在传播早期充当信息放大器，表明异构主体的共存会产生非平凡的协同传播效应。

随着大语言模型技术的突破性进展\cite{brown2020language}，AI代理正经历从规则驱动向大语言模型驱动的范式转变。Yang和Ferrara\cite{yang2024anatomy}对Fox8 Botnet的解剖分析揭示，基于大语言模型的社交机器人已能生成高度拟人化内容并协同传播特定信息，模糊了人类与机器的行为边界。从认知理论视角来看，两类主体在信息加工机制上存在根本性差异：Petty和Cacioppo\cite{petty1986elaboration}的精细加工可能性模型指出，人类认知加工沿中心路径与边缘路径展开，受动机水平与认知能力调控；Corbetta和Shulman\cite{corbetta2002control}进一步揭示了人类注意系统的目标导向与刺激驱动双网络结构。这些认知机制意味着人类传播行为具有内在的随机性与情境依赖性，而AI代理基于确定性的语言模型生成内容，其行为由模型参数与提示策略决定。这种根本差异预示着两类主体会产生可区分的传播行为模式，但现有研究尚未在统一的评估框架下对此进行系统性的量化分析。

\subsection{现有研究总结}

综合以上两个方向的研究现状，可以发现现有工作在以下三个层面存在不足，构成了本文的研究切入点：

\textbf{（1）消息层评估的主体感知缺失。}现有的评论排序与价值评估方法\cite{hsu2009ranking,glenski2017predicting}面向同质用户群体设计，以互动量、文本特征或用户影响力等指标衡量评论价值，未将发布主体类型作为特征引入分析框架。在多主体共存环境下，人类用户与AI代理因认知机制不同而呈现差异化的内容生成模式与传播价值分布，现有方法无法捕捉这种由主体异质性驱动的差异。

\textbf{（2）级联层预测的交互类型盲区。}现有的级联预测模型\cite{li2017deepcas,cao2017deephawkes,zhou2021casflow}假设传播节点具有同质性，在建模传播路径时未区分节点的主体身份，无法编码路径上相邻节点间的异构交互类型（H$\to$H / H$\to$A / A$\to$H / A$\to$A）。在多主体环境下，不同交互类型对讨论延续的影响可能存在显著差异，忽略这一维度将导致预测模型丢失重要的结构信息。

\textbf{（3）数据支撑层面的标注空白。}广泛使用的社交媒体传播数据集\cite{ma2016detecting,ma2017detect}构建于大语言模型代理大规模参与社交媒体之前，不包含发布主体类型（人类/AI代理）的标注信息，无法支撑多主体传播行为的对比研究与模型训练。

针对上述三方面不足，本文从消息层、级联层和交互类型层三个评估维度展开研究，构建包含主体类型标注的数据集，提出具有多主体感知能力的评估模型，设计并实现完整的多主体信息认知传播评估系统。需要说明的是，本文涉及的深度表格学习方法（TabTransformer、AutoInt）、序列建模方法（LSTM）、预训练语言模型（BERT、RoBERTa）以及Transformer架构等技术基础将在第二章中详细介绍。

\section{本文研究内容}

本文围绕多主体信息认知传播评估这一总目标，以传播行为指标作为认知过程的代理变量，从消息层、级联层和交互类型层三个评估维度展开研究。具体研究内容包括以下三个方面：

（1）基于组感知融合的评论传播价值评估模型。针对消息层评估维度，本文提出 Group-aware Fusion TabTransformer 模型，将主体标识（agent\_type）作为显式特征组引入，并通过门控融合机制计算样本级别的特征组权重，自适应调整外生上下文、文本语义、情感倾向和主体标识四组异构特征的贡献度，实现评论级别的传播价值量化评估。

（2）基于交互感知的异构路径传播延续预测模型。针对级联层与交互类型层评估维度，本文提出交互感知的异构路径编码器（IHPE）模型，以路径序列为建模单元，通过交互类型编码层显式建模路径上四种主体交互模式（H$\to$H / H$\to$A / A$\to$H / A$\to$A），捕捉从根节点到当前节点的传播上下文依赖，实现路径级别的传播延续预测。

（3）多主体信息认知传播评估系统的设计与实现。基于上述两个核心算法模块，本文设计并实现了完整的评估系统，涵盖数据导入与预处理、评论传播价值评估、传播延续预测、多主体对比分析以及评估结果展示与管理等功能，将两个模型的输出按主体类型聚合，从三个评估维度系统性地量化人类与AI代理的传播行为差异。

\section{论文组织结构}

本文共分为六章，各章内容安排如下：

第一章为绪论，阐述研究背景与意义，综述国内外研究现状，明确本文研究内容与论文组织结构。

第二章为相关工作，介绍本文涉及的相关理论与技术基础。

第三章为系统需求分析与概要设计，围绕多主体信息认知传播评估系统展开需求分析，从功能性需求与非功能性需求两个方面梳理系统应具备的能力，并给出系统的概要设计方案，包括架构设计、模块设计、流程设计与数据库设计。

第四章为基于组感知融合的评论传播价值评估模型，对应消息层评估维度，阐述模型的问题分析、任务定义与模型结构设计，并通过对比实验与消融实验验证模型的有效性。

第五章为基于交互感知的异构路径传播延续预测模型，对应级联层与交互类型层评估维度，阐述模型的问题分析、任务定义与方法设计，并通过对比实验与消融实验验证模型的有效性。

第六章为系统实现与测试，详细阐述各功能模块的具体实现，并通过功能测试与性能测试验证系统的正确性与可用性。